% Chapter 8: Words 1751 - 2000
\chapter{Advanced Mastery}
\label{chap:chapter8}

{\small\sffamily\color{mutedtext} Final frequency tier reaching full 2,000-word mastery for advanced reading and comprehension.\par}

\vspace{0.4em}
\markboth{Chapter 8: Advanced Mastery}{Words 1751--2000}

\begin{multicols}{2}
\vocentry{1751}{Краска}{'kraska}{Paint, dye}{Краска на этой чашке уже поцарапалась, но она всё равно моя любимая.}{The paint on this cup is scratched already but still it’s my favorite one.}
\vocentry{1752}{Разница}{'raznitsa}{A difference}{Разница в возрасте никак не влияет на нашу дружбу.}{The age difference doesn’t affect our friendship in any way.}
\vocentry{1753}{Расстояние}{rasta'janije}{A distance}{Это расстояние слишком большое для меня. Я не привыкла столько ходить пешком.}{This distance is too long for me. I’m not used to walking on foot so much.}
\vocentry{1754}{Кружок}{kru'zhok}{An interest group; a circle (small)}{Моя сын посещает художественный кружок после уроков.}{After classes my son attends an art interest group.}
\vocentry{1755}{Пожилой}{pazhi'loj}{Elderly}{Пожилой мужчина попросил меня объяснить ему, как пользоваться банкоматом.}{An elderly man asked me to explain to him how to use a cash machine.}
\vocentry{1756}{Непонятно}{nepa'njatnə}{It is unclear (impersonal); incomprehensively}{Мне непонятно, почему женщинам часто платят меньше, чем мужчинам.}{It is unclear to me why women are often paid less than men.}
\vocentry{1757}{Окончательно}{akan'tchatel’nə}{Finally, definitevely}{Мы опубликуем только окончательно утверждённый текст статьи.}{We’ll publish only the finally approved text of the article.}
\vocentry{1758}{Дон}{don}{The Don River}{Дон протекает в европейской части России.}{The Don River flows through the European part of Russia.}
\vocentry{1759}{Заставлять}{zastav'l’at’}{To make, to force}{Зачем заставлять ребёнка есть суп, если он не хочет?}{Why make the child eat soup if he doesn’t want to?}
\vocentry{1760}{Удивительный}{udi'vitel’nyj}{Amazing, marvelous, wonderful}{Поздравляю тебя с днём рождения! Ты удивительный человек, и я рад, что ты есть в моей жизни.}{Happy birthday to you! You’re an amazing person and I’m glad I have you.}
\vocentry{1761}{Страдать}{stra'dat’}{To suffer}{Посевы продолжают страдать от засухи.}{The crops are continuing to suffer from drought.}
\vocentry{1762}{Седой}{se'doj}{Grey-haired, grey}{Ты всё принимаешь близко к сердцу, поэтому ты уже седой. А ты ещё такой молодой!}{You take everything close to heart which is why you’re grey-haired already. And you’re still so young!}
\vocentry{1763}{Политика}{pa'litika}{A policy; politics}{Новая экономическая политика обещает большие перемены.}{The new economic policy promises great changes.}
\vocentry{1764}{Поколение}{paka'lenije}{A generation}{Каждое новое поколение думает, что оно лучше остальных.}{Every new generation think they’re better than the one before.}
\vocentry{1765}{Ожидание}{azhi'danije}{Waiting, expectation}{Ожидание хуже всего на свете: ты ничего не можешь сделать, и от тебя ничего не зависит.}{Waiting is worse than anything else in the world: you can’t do anything, and nothing depends on you.}
\vocentry{1766}{Записать}{zapi'sat’}{To write down; to register, to enroll}{У тебя есть лист бумаги и ручка? Я хочу записать этот рецепт.}{Have you got a sheet of paper and a pen? I want to write down this recipe.}
\vocentry{1767}{Фонарь}{fa'nar’}{A lamppost; a lantern}{Этот фонарь у моего дома такой яркий! Он мешает мне уснуть по ночам.}{This lamppost by my house is so bright! It prevents me from sleeping at night.}
\vocentry{1768}{Куча}{'kutcha}{A pile, a heap}{Эта куча бумаг на твоём столе с каждым днём становится всё больше и больше!}{This pile of papers on your desk is getting bigger and bigger every day!}
\vocentry{1769}{Полоса}{pala'sa}{A streak, a strip}{Надеюсь, теперь полоса неудач в моей жизни закончилась.}{I hope the streak of bad luck in my life is over now.}
\vocentry{1770}{Страсть}{strast’}{A passion}{Пение это не просто её хобби, это её страсть.}{Singing is not just a hobby of hers, it’s her passion.}
\vocentry{1771}{Сержант}{ser'zhant}{A sergeant (military)}{Этот молодой сержант полиции помогает мне с расследованием.}{This young police sergeant is helping me with the investigation.}
\vocentry{1772}{Святой}{sv’ə'toj}{A saint; saint}{Этот святой считается покровителем торговли.}{This saint is considered to be the patron of trade.}
\vocentry{1773}{Повесть}{'povest’}{A story, a tale (a genre of narrative, something between a short story and a novel in size)}{Эта повесть рассказывает о неразделённой любви.}{This story tells about unanswered love.}
\vocentry{1774}{Шёпот}{'shopət}{A whisper}{Я слышала тихий шёпот в соседней комнате, но не смогла разобрать, кто говорит.}{I heard a quiet whisper in the next room but couldn’t make out who was talking.}
\vocentry{1775}{Теория}{te'orija}{A theory}{Его теория основана на многолетних исследованиях.}{His theory is based on longstanding research.}
\vocentry{1776}{Интересовать}{interesa'vat’}{To interest, to concern}{Почему меня должна интересовать твоя личная жизнь?}{Why should your personal life interest me?}
\vocentry{1777}{Чисто}{tchis'tə}{Clean (adv or impersonal), neatly}{Дома так чисто! Когда вы успели всё убрать?}{The house is so clean! When did you have time to tidy everything up?}
\vocentry{1778}{Пояс}{'pojas}{A belt}{Тонкий кожаный пояс красиво подчёркивал её талию.}{A thin leather belt circled her waist beautifully.}
\vocentry{1779}{Близко}{'blizkə}{Close (adv) (both figuratively and literally), near}{Мы жили близко друг к другу и провели всё детство вместе.}{We used to live close to each other and spent our whole childhood together.}
\vocentry{1780}{Включить}{vklu'tchit’}{To turn on; to include (perfective)}{Уже пора включить свет. Я не вижу, что читаю.}{It’s time to turn on the light. I don’t see what I’m reading.}
\vocentry{1781}{Поразить}{para'zit’}{To strike, to amaze; to strike, to hit}{Её талант просто не мог не поразить жюри.}{Her talent just couldn’t but strike the jury.}
\vocentry{1782}{Вокзал}{vak'zal’}{A railway station}{Я отвезу тебя на вокзал и подожду, пока не прибудет твой поезд.}{I’ll take you to the railway station and wait until your train arrives.}
\vocentry{1783}{Вещество}{vestchest'vo}{A substance, a matter}{Мы отправили вещество в лабораторию для более подробного анализа.}{We sent the substance to the laboratory for a more detailed analysis.}
\vocentry{1784}{Частный}{'tchasnyj}{Private}{Частный вертолёт президента приземлился за зданием парламента.}{The president’s private helicopter landed behind the Parliament building.}
\vocentry{1785}{Корень}{'koren’}{A root}{Корень этого растения целебный. Его можно заваривать и пить как чай.}{The root of this plant is healing. It can be infused and drunk as tea.}
\vocentry{1786}{Помолчать}{pamal'tchat’}{To keep silent, to be silent (for a while)}{Дети, неужели вы не можете помолчать хотя бы минуту?}{Kids, can’t you really keep silent for at least a minute?}
\vocentry{1787}{Столовый}{sta'lovyj}{Table (adj)}{Этот столовый сервиз сделан из фарфора.}{This table set is made of china.}
\vocentry{1788}{Ткань}{tkan’}{A cloth}{Эта гладкая шёлковая ткань подойдёт для платья, которое я шью.}{This smooth silk cloth will suit the dress I’m sewing.}
\vocentry{1789}{Открыться}{at'krytsa}{To open (reflexive, perfective); to confide}{Новый аквапарк должен открыться через неделю.}{The new water park must open in a week.}
\vocentry{1790}{Поток}{pa'tok}{A torrent, a stream}{Горный поток заканчивался высоким водопадом.}{The mountain torrent ended with a high waterfall.}
\vocentry{1791}{Вынужденный}{'vynuzhdennyj}{Forced}{Вынужденный отъезд отца расстроил всю семью.}{Father’s forced departure upset the whole family.}
\vocentry{1792}{Полиция}{pa'litsija}{Police}{Полиция не смогла предотвратить преступление.}{The police was unable to avert the crime.}
\vocentry{1793}{Стыдный}{'stydnyj}{Not used in the infinitive, the most widespread form is the impersonal one - стыдно {\ipafont [stydn{\ipafont ə}]} - To be ashamed}{Я не понимаю, почему мне должно быть стыдно за то, что я сказала правду.}{I don’t understand why I should be ashamed of having told the truth.}
\vocentry{1794}{База}{'baza}{A base}{Военная база, на которой работает мой отец, – секретный объект.}{The military base my father works at is a secret location.}
\vocentry{1795}{Слышно}{'slyshnə}{One can hear (impersonal)}{Отсюда мне хорошо слышно, что они говорят. Я могу разобрать каждое слово.}{From here I can hear very well what they’re saying. I can make out every single word.}
\vocentry{1796}{Строгий}{'strogij}{Strict}{Её строгий вид – всего лишь маска. На самом деле она добрая и понимающая.}{Her strict looks are just a mask. In reality she’s kind and understanding.}
\vocentry{1797}{Шинель}{shi'nel’}{A uniform overcoat (military)}{Шинель солдата промокла насквозь, и он боялся что может простудиться.}{The soldier’s uniform overcoat got wet through, and he was afraid he might catch a cold.}
\vocentry{1798}{Блестящий}{bles't’astchij}{Brilliant; shining, glittering}{Твой дядя блестящий адвокат. Я впечатлен его знаниями!}{Your uncle is a brilliant lawyer. I’m impressed by his knowledge!}
\vocentry{1799}{Выбор}{'vybor}{A choice}{Я думаю, ты сделала правильный выбор. Семья важнее, чем деньги.}{I think you’ve made the right choice. Family is more important than money.}
\vocentry{1800}{Сведения}{'svedenija}{Information, intelligence}{Эти сведения из надёжного источника. Не сомневайся.}{This information is from a reliable source. Have no doubts.}
\vocentry{1801}{Комиссия}{ka'misija}{A commission, a committee}{Для проверки состояния завода была создана специальная комиссия.}{A special commission was created to check the condition of the factory.}
\vocentry{1802}{Направить}{nap'ravit’}{To direct}{Мы решили направить ваше дело в апелляционный суд.}{We’ve decided to direct your case to the court of appeals.}
\vocentry{1803}{Палка}{'palka}{A stick, a cane}{Мальчик притворился, что палка – это его винтовка.}{The boy pretended his stick was his rifle.}
\vocentry{1804}{Запомнить}{za'pomnit’}{To memorize}{Я не могу запомнить её имя и не знаю, как к ней обращаться.}{I can’t memorize her name and don’t know how to address her.}
\vocentry{1805}{Бороться}{ba'rotsa}{To fight, to struggle}{Тебе нужно бороться с твоей депрессией. Так не может больше продолжаться.}{You need to fight with your depression. It can’t go on like that anymore.}
\vocentry{1806}{Везти}{ves'ti}{To drive (to carry by vehicle); to be lucky}{Мне пришлось самой везти детей в клинику, потому что муж был на работе.}{I had to drive the children to the clinic myself because my husband was at work.}
\vocentry{1807}{Колонна}{ka'lona}{A column, a pillar}{Высокая колонна в центре площади была воздвигнута в честь великой победы.}{The tall column in the middle of the square was erected in honor of the great victory.}
\vocentry{1808}{Лёд}{ljod}{Ice}{У нас есть лёд в холодильнике? Я хочу охладить свой напиток.}{Do we have any ice in the fridge? I want to cool my drink.}
\vocentry{1809}{Замечательный}{zame'tchatel’nyj}{Remarkable, outstanding}{В той ситуации ты проявил замечательный профессионализм.}{You showed remarkable professionalism in that situation.}
\vocentry{1810}{Надпись}{'natpis’}{An inscription}{Странная надпись на стене пещеры привлекла внимание многих учёных.}{A strange inscription on the cave wall attracted the attention of many scientists.}
\vocentry{1811}{Одинокий}{adi'nokij}{Lonely, lonesome}{У меня нет семьи, но я не могу сказать, что я одинокий человек.}{I don’t have a family but I can’t say I’m a lonely person.}
\vocentry{1812}{Пиджак}{pid'zhak}{A jacket (a suit jacket for men)}{На улице было холодно, и он любезно предложил ей свой пиджак.}{It was cold outside and he kindly offered his jacket to her.}
\vocentry{1813}{Тарелка}{ta'relka}{A plate}{Фарфоровая тарелка выскользнула из её рук и разбилась на сотни кусочков.}{The china plate slipped out of her hands and broke into hundreds of pieces.}
\vocentry{1814}{Научиться}{nau'tchitsa}{To learn}{Я хочу научиться играть на музыкальном инструменте. Например, на скрипке.}{I want to learn how to play a musical instrument. The violin, for example.}
\vocentry{1815}{Пища}{'pistcha}{Food}{Пища должна быть здоровой и сбалансированной.}{Food must be healthy and balanced.}
\vocentry{1816}{Назначить}{naz'natchit’}{To appoint, to designate}{Во вторник президент собирается назначить новых министров.}{On Tuesday the president is going to appoint new ministers.}
\vocentry{1817}{Подруга}{pad'ruga}{A friend (female)}{Моя лучшая подруга всегда знает, что мне нужно.}{My best friend always knows what I need.}
\vocentry{1818}{Выбирать}{vybi'rat’}{To choose, to select}{Она совсем не умеет выбирать одежду и часто выглядит нелепо.}{She’s absolutely unable to choose clothes and often looks ridiculous.}
\vocentry{1819}{Взрыв}{vzryv}{An explosion}{Взрыв причинил огромный ущерб машинам, припаркованным перед зданием.}{The explosion did enormous harm to the cars parked in front of the building.}
\vocentry{1820}{Сущность}{'sustchnəst’}{Essence, nature}{В сложных обстоятельствах люди проявляют свою истинную сущность.}{In hard circumstances people show their true essence.}
\vocentry{1821}{Батарея}{bata'reja}{A battery (electrical); a radiator}{Она разговаривала по телефону всю ночь, пока не села батарея.}{She was talking on the phone the whole night until the battery was out.}
\vocentry{1822}{Техника}{'tehnika}{An equipment, machinery; a technique}{Это сложная техника. Боюсь, я не смогу её починить.}{It’s a complicated piece of equipment. I’m afraid I won’t be able to fix it.}
\vocentry{1823}{Пастух}{pas'tuh}{A shepherd}{Пастух уснул, и стадо разбежалось по всему полю.}{The shepherd fell asleep and the herd scattered about the field.}
\vocentry{1824}{Напомнить}{na'pomnit’}{To remind}{Не забудь напомнить мне поставить будильник.}{Don’t forget to remind me to set the alarm clock.}
\vocentry{1825}{Сомневаться}{samnə'vatsa}{To doubt}{Не позволяй никому заставить тебя сомневаться в правильности своего решения.}{Don’t let anyone make you doubt the correctness of your decision.}
\vocentry{1826}{Крыльцо}{kryl’'tso}{A porch}{Будь осторожен! Крыльцо мокрое после дождя.}{Be careful! The porch is wet after the rain.}
\vocentry{1827}{Радоваться}{'radavatsa}{To rejoice, to be happy}{Способность радоваться за других людей – редкое явление в современном мире.}{The ability to rejoice for other people is a rare phenomenon in the modern world.}
\vocentry{1828}{Покрыть}{pak'ryt’}{To cover (perfective) (both literally and figuratively)}{Этой суммы не достаточно, чтобы покрыть все мои расходы.}{This sum is not enough to cover all my expenses.}
\vocentry{1829}{Глянуть}{'gl’anut’}{To take a look, to glance (to have a brief look, often to check something)}{Я подготовила все документы. Можешь глянуть.}{I’ve prepared all the documents. You may take a look.}
\vocentry{1830}{Значительный}{zna'tchitel’nyj}{Signififcant, considerable}{Эта организация внесла значительный вклад в культурное развитие страны.}{This organisation made a significant contribution to the cultural development of the country.}
\vocentry{1831}{Захватить}{zahva'tit’}{To capture, to seize (perfective)}{Во многих фильмах злодеи пытаются захватить весь мир.}{In many movies villains are trying to capture the whole world.}
\vocentry{1832}{Вождь}{vosht’}{A chief (of a tribe), (often related to Soviet leaders)}{Люди говорили, что Ленин – вождь революции.}{People used to say that Lenin was the chief of the revolution.}
\vocentry{1833}{Восточный}{vas'tochnyj}{Eastern}{Восточный мир полон особенных традиций, которые непонятны большинству европейцев.}{The eastern world is full of special traditions that are unclear to most europeans.}
\vocentry{1834}{Приносить}{prina'sit’}{To bring, to fetch}{Пока ты болеешь, я могу приносить тебе продукты, если хочешь.}{While you’re ill I can bring you some food, if you like.}
\vocentry{1835}{Международный}{mezhduna'rodnyj}{International}{К счастью, международный конфликт был разрешен быстрее, чем ожидалось.}{Luckily, the international conflict was resolved faster than expected.}
\vocentry{1836}{Необходимость}{neapha'diməst’}{A necessity, a need}{Государство признаёт необходимость реформирования политической системы.}{The state admits the necessity of reforming the political system.}
\vocentry{1837}{Оттого}{atta'vo}{That is why}{Он первый раз в таком большом городе, оттого он немного смущён.}{It’s his first time in such a big city, that is why he’s a bit embarrassed.}
\vocentry{1838}{Духовный}{du'hovnyj}{Spiritual}{Этот священник – глубоко духовный человек, почти святой.}{This priest is a deeply spiritual man, almost a saint.}
\vocentry{1839}{Голод}{'goləd}{Hunger}{Голод в оккупированном городе заставлял людей забывать о своих ценностях.}{The hunger in the occupied city made people forget about their values.}
\vocentry{1840}{Девица}{de'vitsa}{A damsel, a maid}{Девица, попавшая в беду, – классический сюжет многих книг и фильмов.}{A damsel in distress is the classical plot of many books and movies.}
\vocentry{1841}{Временный}{'vremennyj}{Temporary}{Говорят, что это ограничение носит временный характер, но я не верю этому.}{They say this restriction is of a temporary nature, but I don’t believe it.}
\vocentry{1842}{Составить}{sas'tavit’}{To constitute, to comprise; to construct (perfective)}{Местные компании должны составить ядро будущей организации.}{The local companies must constitute the core of the future organisation.}
\vocentry{1843}{Ровно}{'rovnə}{Sharp, exactly; smoothly, evenly}{Я заеду за тобой ровно в шесть. Будь готова к этому времени.}{I’ll pick you up at six sharp. Be ready by this time.}
\vocentry{1844}{Мечта}{metch'ta}{A dream}{Дом у океана – это моя мечта. Я делаю всё возможное, чтобы её осуществить.}{A house by the ocean is my dream. I’m doing everything possible to make it come true.}
\vocentry{1845}{Кольцо}{kal’'tso}{A ring}{Это золотое кольцо дорого мне как память о бабушке.}{This golden ring is dear to me as a memory of my grandmother.}
\vocentry{1846}{Ага}{a'ga}{Yep, yeah (colloquial for ‘yes’)}{Ага, хорошо. Обязательно зайду к тебе, когда будет время.}{Yep, alright. I’ll surely drop by if I have time.}
\vocentry{1847}{Дочка}{'dotchka}{A daughter (diminutive or colloquial for ‘дочь’ {\ipafont [dotch’]})}{Я не могу поверить, что наша маленькая дочка уже выросла.}{I can’t believe our little daughter has grown up already.}
\vocentry{1848}{Королева}{kara'leva}{A queen}{Английская королева обладает только номинальной властью.}{The English Queen has only nominal power.}
\vocentry{1849}{Картошка}{kar'toshka}{Potatoes (colloquial)}{На ужин у нас жареная картошка и овощной салат.}{We have fried potatoes and a vegetable salad for supper.}
\vocentry{1850}{Достоинство}{das'toinstvə}{Dignity}{Я не собираюсь терять своё достоинство и умолять их взять меня обратно на работу.}{I’m not going to lose my dignity and beg them to take me back to work.}
\vocentry{1851}{Медведь}{med'ved’}{A bear}{Медведь – популярный персонаж русских сказок.}{The bear is a popular character of Russian fairy tales.}
\vocentry{1852}{Содержание}{sader'zhanije}{Content; maintenance}{Меня интересует содержание этого курса. Может, стоит на него подписаться?}{I’m interested in the content of this course. Maybe it’s worth signing up for it.}
\vocentry{1853}{Проект}{pra'ekt}{A project}{Можно мне взглянуть на проект здания? Я хочу убедиться, что всё в порядке.}{May I have a look at the project of the building? I want to make sure everything is alright.}
\vocentry{1854}{Обида}{a'bida}{A grudge, an insult}{Твой отказ приехать – это личная обида для меня.}{Your refusal to come is a personal insult for me.}
\vocentry{1855}{Согласен}{sag'lasen}{To agree (part of a predicative construction, the initial form is for a masculine or neuter subject)}{Он согласен со всем, что говорят другие. У него нет своего мнения.}{He agrees with everything others say. He doesn’t have his own opinion.}
\vocentry{1856}{Покупать}{paku'pat’}{To buy}{Я люблю покупать вещи онлайн. Никакой суеты и тяжёлых сумок.}{I like to buy things online. No fuss and heavy bags.}
\vocentry{1857}{Сперва}{sper'va}{First, at first, firstly (more typical for literary style)}{Сперва помой руки, а потом садись за стол.}{Wash your hands first and then sit at the table.}
\vocentry{1858}{Юный}{'junyj}{Young, youthful (in especially young, teenage years)}{Он такой юный и наивный и совсем не знает жизни.}{He’s so young and naive and doesn’t know life at all.}
\vocentry{1859}{Любопытство}{l’uba'pytstva}{Curiosity}{Её любопытство раздражает меня! Она хочет знать всё обо всех в мельчайших деталях.}{Her curiosity irritates me! She wants to know everything about everybody in smallest detail.}
\vocentry{1860}{Создание}{saz'danije}{Creation; a creature}{Создание такой информационной базы может занять годы.}{The creation of such a database may take years.}
\vocentry{1861}{Паспорт}{'paspart}{A passport}{Только не говори мне, что ты забыла свой паспорт дома. Тебя не пустят на борт.}{Just don’t tell me you left your passport at home. They won’t let you on board.}
\vocentry{1862}{Двести}{'dvesti}{Two hundred}{Мне предложили на двести долларов больше, если я закончу работу раньше.}{I was offered an extra two hundred dollars if I finish the work earlier.}
\vocentry{1863}{Коллега}{ka'lega}{A colleague}{Мой коллега на больничном, и мне приходится работать за двоих.}{My colleague is on a sick leave and I have to work for both of us.}
\vocentry{1864}{Девчонка}{dev'tchjonka}{A girl, a gal (colloquial)}{Моя племянница – милая, активная девчонка.}{My niece is a nice and active girl.}
\vocentry{1865}{Луч}{lutch’}{A ray, a beam}{Котёнок бегал по комнате, пытаясь поймать луч света.}{A kitten was running about the room trying to catch a light ray.}
\vocentry{1866}{Стучать}{stu'tchat’}{To knock}{Тебя не учили стучать в дверь? Это моя комната, моё пространство.}{Haven’t you been taught how to knock on the door? It’s my room, my space.}
\vocentry{1867}{Коммунист}{kamu'nist}{A communist (person who follows a communist or Marxist-Leninist philosophy)}{Мой дедушка – убеждённый коммунист. Спорить с ним бесполезно.}{My granddad is a convinced communist. It’s no use arguing with him.}
\vocentry{1868}{Просьба}{'pros’ba}{A request}{Твоя просьба неуместна. Я не могу помогать тебе бесконечно.}{Your request is inappropriate. I can’t help you endlessly.}
\vocentry{1869}{Домик}{'domik}{A cottage, a little house}{У моих друзей есть домик в лесу. Идеальное место, чтобы отдохнуть от суеты города.}{My friends have a cottage in the forest. A perfect place to rest from the fuss of the city.}
\vocentry{1870}{Полчаса}{poltcha'sa}{Half an hour}{Я немного занята сейчас. Перезвоню тебе через полчаса.}{I’m a bit busy now. I’ll call you back in half an hour.}
\vocentry{1871}{Железо}{zhe'leza}{Iron}{Железо – один из важнейших элементов для правильного обмена веществ.}{Iron is one of the most important elements for a proper metabolism.}
\vocentry{1872}{Отличие}{at'litchije}{A difference; a distinction}{Смотри, какая интересная головоломка. Найди одно отличие между этими картинками.}{Look what an interesting puzzle. Find one difference between these pictures.}
\vocentry{1873}{Менять}{men’'at’}{To change; to replace}{Я ничего не хочу менять в своей жизни. Я абсолютно счастлив.}{I don’t want to change anything in my life. I’m absolutely happy.}
\vocentry{1874}{Мгновенно}{mgna'venna}{Instantly, immediately}{При ограблении сигнализация сработала мгновенно.}{During the robbery the alarm went off instantly.}
\vocentry{1875}{Пёс}{p’os}{A dog, a hound (male)}{Мой преданный пёс всегда встречает меня радостным лаем, когда я возвращаюсь с работы.}{My loyal dog always greets me with joyful barking when I’m returning home from work.}
\vocentry{1876}{Разглядывать}{raz'gljadyvat’}{To look at, to examine (to look at attentively trying to find something or just paying attention to detail)}{Мне нравится разглядывать насекомых под лупой.}{I like to look at insects under the magnifying glass.}
\vocentry{1877}{Горький}{'gor’kij}{Bitter (both figuratively and literally)}{Я не понимаю, как людям может нравится горький шоколад.}{I don’t understand how people can like bitter chocolate.}
\vocentry{1878}{Орудие}{a'rudije}{A tool (both literally and figuratively)}{Ты не понимаешь: ты всего лишь орудие в её хитрой игре.}{You don’t understand: you’re just a tool in her cunning game.}
\vocentry{1879}{Специально}{spetsi'al’nə}{Specially; on purpose}{Ты знаешь, я не люблю виноград. Я купила его специально для тебя.}{You know, I don’t like grapes. I bought them specially for you.}
\vocentry{1880}{Сказка}{'skaska}{A fairy tale, a tale}{Когда я была маленькой девочкой, мне очень нравилась сказка про Белоснежку.}{When I was a little girl I liked the fairy tale about Snow White a lot.}
\vocentry{1881}{Закурить}{zaku'rit’}{To smoke (to start smoking a cigarette) (perfective)}{Можно мне закурить прямо здесь? Или мне лучше выйти на улицу?}{May I smoke right here? Or should I better go outside?}
\vocentry{1882}{Семейный}{se'mejnyj}{Family (adj)}{В субботу у нас будет большой семейный ужин, так я что не пойду с вами.}{On Saturday we’re having a big family supper, so I’m not going with you.}
\vocentry{1883}{Устать}{us'tat’}{To get tired (perfective)}{Как ты мог так быстро устать? Мы прошли всего пару километров.}{How could you get tired so soon? We’ve just walked a few kilometers.}
\vocentry{1884}{Прикрыть}{prik'ryt’}{To close, to shut (softly and not completely); to cover up}{Я уложила ребёнка спать, но забыла прикрыть дверь, и шум голосов разбудил его.}{I put the baby to sleep but forgot to close the door and the noise of voices woke him up.}
\vocentry{1885}{Твёрдо}{'tv’ordə}{Firmly (both literally and figuratively, decisively, resolutely)}{Я твёрдо верю в то, что все перемены в жизни к лучшему.}{I firmly believe that all the changes in life are for the better.}
\vocentry{1886}{Молодец}{mala'dets}{Well done}{Отличный бросок! Молодец!}{Nice throw! Well done!}
\vocentry{1887}{Замолчать}{zamal'tchat’}{To become/fall silent (to stop talking) (perfective)}{Она отвергла все мои аргументы, и я решила, что мне лучше замолчать.}{She rejected all my arguments and I decided I’d better become silent.}
\vocentry{1888}{Завести}{zaves'ti}{To have (to start having, like a family, child, a pet) (perfective)}{Я хочу завести домашнего питомца, но не могу выбрать между кошкой и собакой.}{I want to have a pet but can’t choose between a cat and a dog.}
\vocentry{1889}{Прибыть}{pri'byt’}{To arrive}{Поезд должен был прибыть полчаса назад. Что могло случиться?}{The train was to arrive half an hour ago. What could have happened?}
\vocentry{1890}{Собрание}{sa'branije}{A meeting, a gathering; a collection}{Обязательно посети завтрашнее собрание. Там будут обсуждаться важные вопросы.}{Be sure to attend tomorrow’s meeting. They are going to discuss important matters there.}
\vocentry{1891}{Девять}{'dev’at’}{Nine}{В этом сервизе десять тарелок. Я вижу только девять.}{There are ten plates in this dining set. I see only nine.}
\vocentry{1892}{Слон}{slon}{An elephant}{Слон обладает феноменальной памятью.}{The elephant has a phenomenal memory.}
\vocentry{1893}{Подтвердить}{pattver'dit’}{To confirm}{Вам нужно подтвердить это действие, прежде чем вы перейдёте к следующему.}{You need to confirm this action before you go over to the next one.}
\vocentry{1894}{Юноша}{'junasha}{A youth (male), a young man}{Этот юноша талантлив, но недостаточно трудолюбив.}{This youth is talented but not hard working enough.}
\vocentry{1895}{Сойти}{saj'ti}{To go off; to go down (perfective)}{Дай я помогу тебе сойти с автобуса. Здесь рядом большая лужа.}{Let me help you get off the bus. There’s a big puddle near here.}
\vocentry{1896}{Поступать}{pastu'pat’}{To act, to behave; to enter, to join}{Ты уже не ребёнок. Тебе нужно научиться поступать мудрее.}{You’re not a child anymore. You must learn to act more wisely.}
\vocentry{1897}{Крутой}{kru'toj}{Steep; cool, awesome (colloquial)}{Там впереди крутой спуск. Я думаю, нам лучше сойти с велосипедов и пройти пешком.}{There’s a steep slope ahead. I think we’d better get off the bikes and go on foot.}
\vocentry{1898}{Валить}{va'lit’}{To blame (to put the blame on somebody else)}{Не надо валить всё на меня! Ты тоже был там!}{Don’t blame everything on me! You were there too!}
\vocentry{1899}{Объект}{ab'jekt}{An object}{Вчера я видел в небе странный объект. Надеюсь, это были инопланетяне.}{Yesterday I saw a strange object in the sky. I hope those were extraterrestrials.}
\vocentry{1900}{Убитый}{u'bityj}{Murdered, killed}{Мужчина, убитый в собственном доме, был найден только через три дня.}{The man, murdered in his own house, was found only after three days.}
\vocentry{1901}{Сравнение}{srav'nenije}{A comparison}{Сравнение с животными или предметами иногда помогает писателям лучше описать героя.}{A comparison with animals or objects sometimes helps writers to describe a character better.}
\vocentry{1902}{Справа}{'sprava}{On the right, to the right}{Пожалуйста, принеси лимонад. Он возле микроволновки, справа.}{Please, fetch the lemonade. It’s next to the microwave, on the right.}
\vocentry{1903}{Благородный}{blaga'rodnyj}{Noble}{Ты правильно сделал, что не взял деньги обратно. Это был благородный поступок.}{You did right not taking the money back. That was a noble deed.}
\vocentry{1904}{Океан}{ake'an}{An ocean}{Тихий океан самый большой и самый глубокий океан в мире.}{The Pacific Ocean is the largest and the deepest ocean in the world.}
\vocentry{1905}{Посёлок}{pa's’olək}{A village (usually a big one); a settlement}{К утру весь посёлок знал о вчерашней драке.}{By the morning the whole village knew about yesterday’s fight.}
\vocentry{1906}{Тайга}{taj'ga}{Taiga}{Тайга – это хвойные леса, которые простираются на тысячи километров.}{Taiga are coniferous forests that stretch for thousands of kilometers.}
\vocentry{1907}{Борода}{bara'da}{A beard}{Мне нравится моя борода. Мне кажется, она придаёт мне брутальный вид.}{I like my beard. It seems to me it makes me look more masculine.}
\vocentry{1908}{Пачка}{'patchka}{A pack, a package}{Эта пачка сигарет была в твоём кармане. Ты куришь?}{This pack of cigarettes was in your pocket. Do you smoke?}
\vocentry{1909}{Надоесть}{nada'jest’}{To get tired of, to get bored by (impersonal, the person who’s tired/bored takes the dative case)}{Как тебе могла надоесть эта песня? Я готова слушать её снова и снова.}{How could you get tired of this song? I’m ready to listen to it over and over again.}
\vocentry{1910}{Привычка}{pri'vytchka}{A habit}{Я знаю, что есть по ночам, – это плохая привычка, но не могу от неё избавиться.}{I know that eating at night is a bad habit but I can’t get rid of it.}
\vocentry{1911}{Сумма}{'suma}{A sum}{Вчера из местного банка была украдена крупная сумма денег.}{A large sum of money was stolen from the local bank yesterday.}
\vocentry{1912}{Вывести}{'vyvesti}{To lead (out of); to withdraw}{Девочка смогла вывести младшего брата из горящего дома.}{The girl managed to lead her younger brother out of the burning house.}
\vocentry{1913}{Прыгать}{'prygat’}{To jump, to leap}{Ты будешь прыгать от радости, когда увидишь наш новый дом.}{You’re going to jump with joy when you see our new house.}
\vocentry{1914}{Создавать}{sazda'vat’}{To create}{Мы будем создавать все необходимые условия для улучшения качества образования.}{We’re going to create all the necessary conditions to improve the quality of education.}
\vocentry{1915}{Грубый}{'grubyj}{Rude; rough}{Я не понимаю, почему он такой грубый сегодня. Обычно он очень вежливый.}{I don’t understand why he’s so rude today. Usually he’s very polite.}
\vocentry{1916}{Поездка}{pa'jestka}{A trip}{Это будет замечательная поездка! Мы посетим две страны и пять городов.}{It’s going to be a wonderful trip! We’ll visit two countries and five cities.}
\vocentry{1917}{Убедиться}{ube'ditsa}{To make sure (perfective)}{Подожди минуточку. Я хочу убедиться, что выключила утюг.}{Wait a minute. I want to make sure I’ve switched off the iron.}
\vocentry{1918}{Ненавидеть}{nena'videt’}{To hate}{Я не понимаю, как можно ненавидеть человека, которого ты когда-то любил.}{I don’t understand how it’s possible to hate a person you used to love.}
\vocentry{1919}{Заявление}{zajav'lenije}{An application; a statement}{Мы рассмотрим ваше заявление в течение этой недели.}{We’ll consider your application within this week.}
\vocentry{1920}{Поражение}{para'zhenie}{A defeat}{Он не может пережить своё поражение на шахматном турнире.}{He can’t live over his defeat at the chess tournament.}
\vocentry{1921}{Август}{'avgust}{August}{Август мой самый нелюбимый месяц. Он напоминает мне, что скоро осень.}{August is my least favorite month. It reminds me that soon there will be autumn.}
\vocentry{1922}{Привычный}{pri'vytchnyj}{Habitual, customary}{Я не люблю, когда что-нибудь вмешивается в мой привычный ритм жизни.}{I don’t like it when something interferes with my habitual pace of life.}
\vocentry{1923}{Официально}{afitsi'alnə}{Officially}{Поздравьте нас! Теперь мы официально женаты.}{Congratulate us! We’re officially married now.}
\vocentry{1924}{Передний}{pe'rednij}{Front, fore}{Просторный передний двор – одно из главных преимуществ этого дома.}{A spacious front yard is one of the major advantages of this house.}
\vocentry{1925}{Север}{'sever}{North}{Я думаю, что если мы пойдём на север, то скоро выйдем из леса.}{I think if we go north, we’ll go out of the forest soon.}
\vocentry{1926}{Шляпа}{'shl’apa}{A hat (with brims)}{Эта шляпа отлично дополнит твой деловой образ.}{This hat will greatly compliment your business look.}
\vocentry{1927}{Изо}{'izə}{From, out of (a form of ‘из’ used before consonant clusters)}{На карнавал в Рио приезжают люди изо всех уголков мира.}{People from all the corners of the world come to the Rio Carnival.}
\vocentry{1928}{Молчание}{mal'tchanije}{Silence (the quiet because of the absence of speech)}{Молчание становилось неловким, но никто не знал, как прервать его.}{The silence was getting awkward but nobody knew how to break it.}
\vocentry{1929}{Исчезать}{istche'zat’}{To disappear, to vanish}{Нельзя вот так исчезать, не предупредив никого. Мы волновались.}{You can’t disappear like that without warning anyone. We were worried.}
\vocentry{1930}{Эй}{ɛj}{Hey}{Эй, ты! Подойди сюда!}{Hey, you there! Come here!}
\vocentry{1931}{Топор}{ta'por}{An axe}{В музее мы увидели боевой топор викингов.}{We saw a battle axe of the Vikings in the museum.}
\vocentry{1932}{Ощущать}{astchust'chat’}{To feel, to sense}{Это удивительно, как я могу ощущать твою поддержку даже на расстоянии.}{It’s amazing how I can feel your support even at a distance.}
\vocentry{1933}{Тащить}{tast'chit’}{To drag, to pull}{Сумка была такой тяжёлой, что я не могла нести её и мне пришлось тащить её.}{The bag was so heavy that I couldn’t carry it and had to drag it.}
\vocentry{1934}{Брюки}{'br’uki}{Trousers (more for formal, business style)}{Дай мне салфетку, пожалуйста. Я пролил сок на брюки.}{Give me a napkin, please. I’ve spilled some juice on my trousers.}
\vocentry{1935}{Соответствующий}{saat'vetstvujustchij}{Corresponding, respective}{Я не думаю, что получу работу. У меня отсутствует соответствующий опыт.}{I don’t think I’ll get the job. I don’t have the corresponding experience.}
\vocentry{1936}{Элемент}{ɛle'ment}{An element, a component; an item}{Этот крошечный элемент оборудования играет очень важную роль в его работе.}{This tiny element of the equipment plays a very important role in its work.}
\vocentry{1937}{Одеяло}{ade'jalə}{A blanket}{Пора доставать тёплое одеяло. Я замёрзла прошлой ночью.}{It’s time to take out a warm blanket. I was freezing last night.}
\vocentry{1938}{Мотор}{ma'tor}{An engine, a motor}{Мотор моей машины издаёт какие-то странные звуки. Мне нужно показать его механику.}{My car’s engine is giving out some strange sounds. I need to take it to a mechanic.}
\vocentry{1939}{Везде}{vez'de}{Everywhere, anywhere}{Я искала твои бумаги везде. Ты уверена, что оставила их у меня дома?}{I’ve been looking for your papers everywhere. Are you sure you left them at my place?}
\vocentry{1940}{Дружба}{'druzhba}{Friendship}{Со временем наша дружба стала ещё сильнее.}{With time our friendship got even stronger.}
\vocentry{1941}{Безопасность}{beza'pasnəst’}{Security, safety}{Государство обязано обеспечивать безопасность своих граждан.}{The state is obliged to ensure the security of its citizens.}
\vocentry{1942}{Скрывать}{skry'vat’}{To hide, to conceal}{Супруги не должны ничего скрывать друг от друга. У них не должно быть никаких секретов.}{Spouses shouldn’t hide anything from each other. They shouldn’t have any secrets.}
\vocentry{1943}{Совершить}{saver'shit’}{To commit, to perform (perfective)}{Я не верю, что он мог совершить это ужасное преступление! Это совсем на него не похоже!}{I don’t believe he could commit that terrible crime! It’s absolutely unlike him.}
\vocentry{1944}{Крест}{krest}{A cross}{Крест – это центральный символ христианской религии.}{The cross is the central symbol of the Christian religion.}
\vocentry{1945}{Земной}{zem'noj}{Earthly, earth (adj), terrestrial}{Наш земной мир не совершенен, но есть надежда, что небесный сильно от него отличается.}{Our earthly world is not perfect but there’s hope that the heavenly one differs from it a lot.}
\vocentry{1946}{Школьный}{'shkol’nyj}{School (adj)}{Я опоздал на школьный автобус, и маме пришлось везти меня в школу на машине.}{I missed the school bus and Mom had to take me to school by car.}
\vocentry{1947}{Замереть}{zame'ret’}{To freeze, to stand motionless (perfective)}{Это простая игра. Как только я скажу волшебные слова, вы должны замереть.}{It’s a simple game. As soon as I say the magic words you must freeze.}
\vocentry{1948}{Прозрачный}{praz'rachnyj}{Transparent, clear}{Бульон для этого супа должен быть прозрачный, а твой мутный и тёмный.}{The broth for this soup must be transparent and yours is cloudy and dark.}
\vocentry{1949}{Подавать}{pada'vat’}{To give (means to give something to someone who’s far from the object)}{Малыш, я больше не буду подавать тебе игрушки, которые ты нарочно сбрасываешь на пол.}{Baby I’m not going to give you the toys your’re dropping on the floor on purpose.}
\vocentry{1950}{Истинный}{'istinyj}{True; genuine}{Он истинный фанат футбола. Он знает всё обо всех командах и их игроках.}{He’s a true football fan. He knows everything about all teams and their players.}
\vocentry{1951}{Кабина}{ka'bina}{A cab, a cabin}{В моём новом грузовике очень удобная кабина: мягкие сиденья, кондиционер.}{My new truck has a very comfortable cab: soft seats, an air conditioner.}
\vocentry{1952}{Пальто}{pal’'to}{A coat, an overcoat}{Я коплю деньги на новое пальто из дизайнерской коллекции.}{I’m saving up money for a new coat from a designer collection.}
\vocentry{1953}{Восторг}{vas'torg}{Delight, ecstacy}{Тебе нужно было видеть его восторг, когда он увидел, что мне удалось купить билеты.}{You should have seen his delight when he saw that I had managed to buy the tickets.}
\vocentry{1954}{Экономика}{ɛka'nomika}{Economy; economics}{Рыночная экономика не может развиваться в жёстких рамках.}{Market economy can’t develop in a rigid framework.}
\vocentry{1955}{Забота}{za'bota}{Care}{Материнская забота – это вещь, которую ничем нельзя заменить.}{Mother’s care can’t be replaced with anything.}
\vocentry{1956}{Произвести}{praizves'ti}{To produce (perfective)}{В этом году наш завод планирует произвести в два раза больше товаров, чем в прошлом году.}{This year our factory is planning to produce twice as many goods as in the previous year.}
\vocentry{1957}{Дождаться}{dazh'datsa}{To wait (to wait for a long time and finally get what you’ve been waiting for), (perfective)}{Я не могу дождаться, когда впервые увижу моего новорождённого племянника.}{I can’t wait to see my newborn nephew for the first time.}
\vocentry{1958}{Талант}{ta'lant}{A talent, a gift}{У него настоящий талант к рисованию. Отправьте его в художественную школу.}{He has a real talent for painting. Send him to an art school.}
\vocentry{1959}{Плыть}{plyt’}{To swim}{Не пытайся плыть против течения. Ты недостаточно силён для этого.}{Don’t try to swim against the current. You’re not strong enough for it.}
\vocentry{1960}{Запись}{'zapis’}{A recording (of a voice or a sound); an entry}{Послушайте эту запись. Вы узнаёте этот голос?}{Listen to this recording. Do you recognize this voice?}
\vocentry{1961}{Ведро}{ved'ro}{A bucket}{Она взяла ведро с водой и швабру и принялась за уборку.}{She took a bucket with water and a mop and got down to work.}
\vocentry{1962}{Крестьянин}{krest’'janin}{A peasant}{Не каждый крестьянин в те времена мог позволить себе такой большой дом.}{Not every peasant in those times could afford such a big house.}
\vocentry{1963}{Поворот}{pava'rot}{A turning; a turn}{Я ехала слишком быстро и пропустила свой поворот.}{I was driving too fast and missed my turning.}
\vocentry{1964}{Подробность}{pad'robnəst’}{A detail, a particularity}{Как прошло твоё свидание? Я хочу знать каждую подробность.}{How did your date go? I want to know every detail.}
\vocentry{1965}{Столб}{stolp}{A post, a pole}{Посмотри, этот заборный столб сгнил. Нам нужно заменить его.}{Look, this fence post is rotten. We need to replace it.}
\vocentry{1966}{Канал}{ka'nal}{A channel; a canal}{Переключи на новостной канал. Мне интересно, что происходит в мире.}{Switch to the news channel. I wonder what’s going on in the world.}
\vocentry{1967}{Пассажир}{pasa'zhyr}{A passenger}{Каждый пассажир самолёта должен иметь действительный паспорт.}{Every passenger of a plane must have a valid passport.}
\vocentry{1968}{Выполнить}{'vypəlnit’}{To carry out, to fulfill (perfective)}{Вы должны выполнить свои обязательства, иначе мы разорвём контракт.}{You must carry out your obligations, otherwise we’ll break the contract.}
\vocentry{1969}{Устройство}{ust'rojstvə}{A device}{Это простое устройство позволяет измерять пульс и кровяное давление.}{This simple device allows you to measure the pulse and the blood pressure.}
\vocentry{1970}{Благодаря}{blagada'r’a}{Thanks to, due to}{Благодаря поддержке родителей и друзей, мне удалось построить впечатляющую карьеру.}{Thanks to the support of my parents and friends I managed to build an impressive career.}
\vocentry{1971}{Направиться}{nap'ravitsa}{To head (reflexive), (perfective)}{Я советую тебе отложить все дела и направиться в горы или к морю.}{I advise you to put away all your business and head to the mountains or to the seaside.}
\vocentry{1972}{Влияние}{vli'janije}{Influence}{Нельзя недооценивать влияние общества на формирование личности ребёнка.}{One can’t underestimate the influence of society on the development of a child’s personality.}
\vocentry{1973}{Ракета}{ra'keta}{A rocket}{Ракета ещё не полностью готова. Нам придётся отложить запуск.}{The rocket is not completely ready yet. We’ll have to delay the launch.}
\vocentry{1974}{Гражданский}{grazh'danskij}{Civil, civic}{Участие в выборах – это мой гражданский долг.}{The participation in the elections is my civil duty.}
\vocentry{1975}{Шутить}{shu'tit’}{To joke; to make fun (of)}{Она совсем не умеет шутить. Мне никогда не нравилось её чувство юмора.}{She’s absolutely unable to joke. I’ve never liked her sense of humor.}
\vocentry{1976}{Смена}{'smena}{A change, replacement; a shift}{Смена правительства должна происходить раз в четыре года.}{The change of government must take place once in four years.}
\vocentry{1977}{Случаться}{slu'tchatsa}{To happen, to occur}{Неудачи должны случаться. Они делают нас сильнее.}{Misfortunes should happen. They make us stronger.}
\vocentry{1978}{Спрятать}{'spr’atat’}{To hide (perfective, transitive)}{Я могу так хорошо спрятать эти деньги, что никто никогда их не найдёт.}{I can hide this money so well that nobody will ever find it.}
\vocentry{1979}{Предстоит}{predsta'it}{Impersonal for ‘To be to’, ‘To have to’, means that the person has something ahead, the subject takes the dative case.}{Я хочу больше узнать об этом городе, раз мне предстоит жить в нём.}{I want to learn more about this city if I am to live in it.}
\vocentry{1980}{Затылок}{za'tylək}{Back of the head, nape}{Он почесал затылок, думая, как ответить на мой вопрос.}{He scratched the back of his head thinking how to answer my question.}
\vocentry{1981}{Юг}{jug}{South}{Летом мы собираемся поехать на юг и отдохнуть у моря.}{In summer we’re planning to go to the south and have a rest at the seaside.}
\vocentry{1982}{Эпоха}{ɛ'poha}{An era, an age, an epoche}{Каждая историческая эпоха имела свои ценности и идеалы.}{Every historical era had its values and ideals.}
\vocentry{1983}{Товар}{ta'var}{Goods (unlike in English has a singular form), a product}{Покупатель заранее оплатил товар и его доставку.}{The buyer paid for the product and its delivery in advance.}
\vocentry{1984}{Проклятый}{'prokl’atyj}{Cursed}{Это проклятый замок, и только настоящая любовь может разрушить заклятье.}{It’s a cursed castle and only true love can break the spell.}
\vocentry{1985}{Ветка}{'vetka}{A twig, a branch; a branch line}{Тонкая ветка ударила меня по лицу и оставила красный след.}{A thin twig hit me in the face and left a red mark.}
\vocentry{1986}{Кормить}{kar'mit’}{To feed}{Животных в зоопарке кормить запрещено.}{It’s forbidden to feed animals in the zoo.}
\vocentry{1987}{Продолжаться}{pradal'zhatsa}{To last, to continue}{Обучение в университете будет продолжаться четыре года.}{The university study will last four years.}
\vocentry{1988}{Разведчик}{raz'vetchik}{An intelligence officer}{Этот разведчик – герой гражданской войны.}{This intelligence officer is a hero of The Civil War.}
\vocentry{1989}{Отлично}{at'litchnə}{Great}{Хорошо, отлично, встретимся у кинотеатра в пять.}{Alright, great, we’ll meet outside the cinema at five.}
\vocentry{1990}{Боже}{'bozhe}{God, Lord}{О, Боже! Я не сохранила документ и придётся начинать сначала.}{Oh, God! I didn’t save the document and I’ll have to start all over again.}
\vocentry{1991}{Молодость}{'moladast’}{Youth}{В этом городе я провёл свою молодость, поэтому я люблю возвращаться сюда.}{I spent my youth in this city and that is why I like to return here.}
\vocentry{1992}{Отдыхать}{atdy'hat’}{To rest, to have rest}{Тебе нужно больше отдыхать. Ты слишком много работаешь.}{You need to rest more. You work too much.}
\vocentry{1993}{Октябрь}{ak't’abr’}{October}{Я точно помню. Это был октябрь.}{I remember it for sure. That was October.}
\vocentry{1994}{Широко}{shira'ko}{Widely, wide (adv)}{Это широко известный производитель кожаной обуви и сумок.}{It’s a widely known producer of leather footwear and bags.}
\vocentry{1995}{Храм}{hram}{A church, a cathedral (more elevated style); a temple}{По воскресеньям мы ходим в храм.}{On Sundays we go to church.}
\vocentry{1996}{Сдать}{zdat’}{To hand in (perfective)}{Учитель сказал ученикам сдать сочинения к четвергу.}{The teacher told the pupils to hand in their essays by Thursday.}
\vocentry{1997}{Шестой}{shes'toj}{Sixth}{Ты ешь уже шестой кекс. Остановись!}{It’s the sixth muffin you’re eating. Stop!}
\vocentry{1998}{Жаловаться}{'zhalavats’a}{To complain}{Я буду жаловаться на ваш сервис!}{I’m going to complain about your service!}
\vocentry{1999}{Скала}{ska'la}{A rock, a cliff}{Нам нужно развернуть корабль. Впереди скала!}{We need to turn the ship around. There is a rock ahead!}
\vocentry{2000}{Расчёт}{ras'tchjot}{Calculation, computation}{Я не смогу произвести такой сложный расчёт без специальной формулы.}{I won’t be able to carry out such a complex calculation without a special formula.}
\end{multicols}
